% Generated draft: requires mathematical and source review.
\documentclass[11pt]{article}
\usepackage[margin=25mm]{geometry}
\usepackage{fontspec}
\setmainfont{Latin Modern Roman}
\usepackage{amsmath,amssymb,mathtools,unicode-math}
\setmathfont{Latin Modern Math}
\usepackage{enumitem,needspace,xurl,hyperref}
\setlength{\parindent}{0pt}
\setlength{\parskip}{5pt}
\setlength{\emergencystretch}{3em}
\newcommand{\N}{\mathbb{N}}
\newcommand{\Z}{\mathbb{Z}}
\newcommand{\Q}{\mathbb{Q}}
\newcommand{\R}{\mathbb{R}}
\newcommand{\C}{\mathbb{C}}
\newcommand{\F}{\mathbb{F}}
\newcommand{\A}{\mathbb{A}}
\newcommand{\PP}{\mathbb{P}}
\newcommand{\E}{\mathbb{E}}
\newcommand{\eps}{\varepsilon}
\newcommand{\dd}{\,\mathrm{d}}
\begin{document}
\section*{UnsolvedMath ID: 30003648}
\textbf{Source problem code:} catalogue entry
\subsection{Definitions and mathematical statement}
For each positive integer \(n\), let \(\mathbb S^n=\{A\in\R^{n\times n}:A^{\mathsf T}=A\}\), with inner product \(\langle M,A\rangle=\operatorname{tr}(MA)\), and put \(\R_+^n=\{x\in\R^n:x_j\geq0\text{ for }1\leq j\leq n\}\). Define the copositive cone and the completely positive cone by
\[
\operatorname{COP}_n=\{M\in\mathbb S^n:x^{\mathsf T}Mx\geq0\text{ for every }x\in\R_+^n\},
\]
\[
\operatorname{CP}_n=\left\{\sum_{i=1}^{r}b_i b_i^{\mathsf T}:r\in\N,\ r\geq1,\ b_i\in\R_+^n\text{ for every }i\right\}.
\]
Zero vectors are allowed among the \(b_i\), so this definition includes the zero matrix.

For a convex cone \(K\subseteq\mathbb S^n\) and a nonzero \(M\in K\), the ray \(\R_+M=\{tM:t\geq0\}\) is an extreme ray if, whenever \(M=U+V\) with \(U,V\in K\), both \(U\) and \(V\) belong to \(\R_+M\). For \(Y\in\mathbb S^n\setminus\{0\}\) and \(\alpha\in\R\), write
\[
H(Y,\alpha)=\{A\in\mathbb S^n:\langle Y,A\rangle=\alpha\}.
\]
This is a supporting hyperplane of \(K\) if \(K\cap H(Y,\alpha)\neq\varnothing\) and, after possibly replacing \((Y,\alpha)\) by \((-Y,-\alpha)\), one has \(\langle Y,A\rangle\geq\alpha\) for every \(A\in K\).

The supplied problem asks for the following descriptions and construction in general dimension:
\begin{enumerate}
\item Give a complete characterization of the nonzero matrices \(M\in\operatorname{COP}_n\) whose rays \(\R_+M\) are extreme, identifying matrices that differ by a positive scalar.
\item Give a complete characterization of the supporting hyperplanes of \(\operatorname{CP}_n\).
\item Given an arbitrary \(A\in\operatorname{CP}_n\), explicitly construct a finite integer \(r\geq1\) and vectors \(b_1,\ldots,b_r\in\R_+^n\) such that
\[
A=\sum_{i=1}^{r}b_i b_i^{\mathsf T}.
\]
\end{enumerate}
The third request concerns explicit construction of a representation whose existence is part of the definition of complete positivity. The source does not specify an input encoding, an exact or approximate computational model, a complexity requirement, a prescribed value of \(r\), or a minimality requirement. No rationality condition on \(A\) or the factors is stated.

The phrase ``in general dimension'' is interpreted here as requesting a description valid for every positive integer \(n\); the source gives no dimension restriction. The required form of a ``complete characterization'' is unspecified. In particular, the source does not explain whether the supporting-hyperplane request concerns all supporting hyperplanes in the definition above, distinguished supporting hyperplanes, or an explicit structural description of their normal matrices. These qualifications remain unresolved.
\subsection{Short English statement}
For arbitrary dimension, describe all extreme rays of the copositive cone and all supporting hyperplanes of the completely positive cone, and construct a nonnegative vector factorization of any given completely positive matrix. The source leaves the required form of the descriptions and the meaning of explicit construction unspecified. Its open status is a supplied dated assessment, not an independently verified conclusion.
\subsection{Sources}
\begin{itemize}
\item Local UnsolvedMath registry, permanent ID \(30003648\), title: \emph{Extremal Copositive Rays and Completely Positive Factorizations}; exact source problem code: \texttt{catalogue entry}. The supplied record belongs to the Oberwolfach Reports Open Problems collection and gives proposed year \(2017\), without naming a proposer.
\item Supplied source citation: \emph{Copositivity and Complete Positivity} (2018), Oberwolfach Reports, pp.\ 3071--3120, DOI \texttt{10.4171/owr/2017/52}. Source URL and supplied literature reference ``Oberwolfach Report 52, copositive and completely positive cones'': \url{https://doi.org/10.4171/owr/2017/52}. The background describes these reports as records of mathematical workshops and their open-problem discussions.
\item Supplied literature reference: Dutour Sikiri\'c--Sch\"urmann--Vallentin, ``rational cp factorizations'': \url{https://doi.org/10.1016/j.laa.2017.02.017}. This is the bibliographic description supplied in the record; no additional bibliographic details or hypotheses are supplied here.
\item Literature assessment date: \texttt{2026-08-21}. The supplied dated curation labels the problem open and states that complete descriptions of the copositive cone are known only in low dimensions, while a general characterization of its extreme rays and unrestricted completely positive factorization remain open. The background identifies this as dated literature triage derived from the supplied OWR curation and calls for independent verification of time-sensitive status and cited claims. This assessment is preserved as supplied dated curation and has not been independently verified.
\item The supplied statement status is \texttt{corrected\_verified}; its accompanying explanation says that the mathematical content was preserved while notation, typography, or omitted definitions were repaired for self-contained publication. This is source metadata, not independent mathematical verification of this draft.
\end{itemize}
\end{document}
